\documentclass[11pt,letterpaper]{article}
% Shared declarations; root documents always import a concrete semantic module.
\usepackage{business-admin}
\usepackage{amsmath,amssymb}
\usepackage{longtable,tabularx,array,booktabs}
\usepackage{enumitem}
\usepackage{ragged2e}
\newcolumntype{Y}{>{\RaggedRight\arraybackslash}X}
\newcommand{\SuiteMeta}[2]{\def\DocTitle{#1}\def\DocSubtitle{#2}\def\DocVersion{1.0}\def\DocStatus{Educational reference}\def\DocOwner{Small Business Management Suite}\def\DocSummary{Original teaching and practice resource; no real-business approval implied.}}
\newcommand{\SuiteOpen}{\makedoctitle\vspace{-0.7em}\noindent\rule{\linewidth}{0.45pt}\par\medskip}
\newcommand{\SourceNote}[2]{\par\medskip\noindent{\footnotesize\textbf{Reading anchor:} Longenecker, Petty, Palich and Hoy, \emph{Small Business Management: Launching \& Growing Entrepreneurial Ventures}, 20th ed., Ch.~#1, textbook p.~#2. This document is independently authored commentary and exercises, not a reproduction of the chapter.}}
\newcommand{\WritingLines}[1]{\par\smallskip\noindent\rule{\linewidth}{0.35pt}\par\vspace{0.35cm}\noindent\rule{\linewidth}{0.35pt}\par\vspace{0.35cm}\noindent\rule{\linewidth}{0.35pt}\par\medskip}
\newcommand{\FormField}[1]{\par\medskip\noindent\textbf{#1}\par\WritingLines{3}}
\newcommand{\SBMCheckBox}{\(\square\)}
\newcommand{\AssessmentRow}[2]{#1 & #2 & \(\square\) & \(\square\) & \(\square\) & \(\square\) & \(\square\)\\\hline}
\newcommand{\AssessmentHeader}{\noindent\begin{tabular}{|p{0.22\linewidth}|p{0.49\linewidth}|c|c|c|c|c|}\hline\textbf{Criterion} & \textbf{Evidence to inspect} & \textbf{0} & \textbf{1} & \textbf{2} & \textbf{3} & \textbf{4}\\\hline}
\newcommand{\AssessmentEnd}{\end{tabular}\par}
\setlength{\emergencystretch}{3em}

\SuiteMeta{Chapter 21: Managing Small Business Operations}{Study guide | entrepreneurship and small business management}
\begin{document}
\SuiteOpen
\section{Learning outcome}
Map an end-to-end operating process and locate its bottleneck.
\section{Prerequisites}
Read chapters 15, 19, 20.\par\textbf{Required baseline:} ability to record assumptions and distinguish observations from opinions.
\section{Chapter concept map}
\begin{itemize}[leftmargin=*,itemsep=3pt]
\item service and manufacturing processes
\item capacity and scheduling
\item inventory records
\item quality methods
\item purchasing and lean production
\end{itemize}

\section{Technical study notes}
\begin{itemize}[leftmargin=*,itemsep=3pt]
\item A process map should distinguish work that changes customer value, inspections, queueing, rework and handoffs so improvement targets the actual constraint.
\item Capacity is limited by bottlenecks and variability; maximizing each station independently can increase inventory and wait time.
\item Quality management includes preventive error-proofing, measured defects, supplier performance and feedback loops rather than final inspection alone.
\end{itemize}

\section{Reading and analysis protocol}
\begin{enumerate}[leftmargin=*,itemsep=5pt]
\item Read the chapter learning objectives, then mark terms and claims that need primary evidence.
\item For each chapter theme, distinguish the entrepreneur decision, information required, parties affected and failure condition.
\item Write a one-paragraph synthesis in your own words; cite the textbook chapter and separate contemporary practice notes.
\item Complete the applied activity and attach the underlying evidence, calculations or interview log.
\end{enumerate}

\section{Worked application and deliverable}
\textbf{Business exercise.} Measure throughput, queue time, failure cost and supplier reliability.\par\textbf{Required submission.} Operations control plan. Include a dated assumption register and an explicit proceed/revise/stop recommendation.\par\textbf{Assurance checks.} Process ownership, defects, capacity and supplier SLAs.
\section{Common reasoning failures}
\begin{itemize}[leftmargin=*,itemsep=3pt]
\item Confusing utilization with throughput
\item Optimizing one workstation at the expense of the process
\end{itemize}
\noindent\textbf{Correction:} require counterevidence and identify the responsible decision-maker before approving the next stage.
\section{Discussion prompt}
Which operational metric would trigger a redesign?\par\medskip\textbf{Evidence expected:} at least two alternatives, one quantitative or independently checkable observation, a trade-off, and a decision threshold.
\section{Exit assessment}
\begin{enumerate}[leftmargin=*,itemsep=5pt]
\item Explain each term from the concept map without copying the textbook.
\item Defend the applied decision using an assumption, evidence and counterargument.
\item Show an objective test or inspection that could prove a key assumption wrong.
\item State what would change the recommendation and who has approval authority.
\end{enumerate}

\SourceNote{21}{567}
\end{document}
